\GXNSFBodyText

\noindent\textbf{项目特色}\par

\noindent\textbf{用完整公开案例贯通方法链}\par
项目以佐佐木希跨越杂志、影视、商业合作和自营视频平台的公开职业轨迹为演示主线，使事件核验、状态识别、扰动分析和跨案例验证围绕同一对象连续展开。案例的价值在于提供可追溯、可比较的长期媒介记录，使机制分析能够落实到连续的证据链中。

\noindent\textbf{拟体现的创新点}\par

\noindent\textbf{平台校准的多模态隐状态模型}\par
在共享条件状态模型之前设置平台专属校准映射，区分“平台指标变了”和“职业状态变了”，解决跨媒介数据量纲、频率和反馈机制不一致的问题。该结构可通过删去校准映射的消融实验直接检验。

\noindent\textbf{具有边界条件的职业韧性测度}\par
把职业韧性拆分为恢复速度、媒介多样性、叙事一致性和波动稳定性，统一归一化方向，并比较“回到原状态”与“转入新稳定状态”两类恢复路径。这一测度的关键在于明确扰动前基线、恢复窗口、各分量的观测依据，以及主分析和敏感性分析的权重口径。

\noindent\textbf{从单案例发现到分层验证的证据链}\par
采用“佐佐木希案例贯通—同构人物样本对照—广西与东盟方法组件测试”的分层路径，据此检验模型和组件分别能够适用到何种情形。该设计将人物职业状态模型的外部效度与跨语言方法组件的适用性分别检验，保证结论与分析单位相匹配。

\noindent\textbf{科学意义}\par

\noindent\textbf{刻画多层媒介网络中的状态迁移}\par
项目把公众人物职业变化转化为多层网络中的动态状态与扰动恢复问题，可补充传统职业阶段理论对平台耦合关系描述不足的缺口。若模型能够在替代样本与不同窗口下保持稳定，将为研究数字平台中的身份持续性、内容迁移和系统韧性提供可计算框架；若不能保持稳定，失效边界本身也构成可检验结论。
